% ch15b Quantum Hall and other topological states, part b (书页 343-352 / 分册p027-036)
\section{劳夫林液体}

在崔琦（Tsui）、斯特默（Störmer）与戈萨德（Gossard）最初的实验（文献 [TSG1982]）中，霍尔电导在 $e^{2}/3h$ 数值处呈现出一个陡峭的平台。在后来的实验中，霍尔电导的分数量化也在 $4/3$、$5/3$、$7/3$、$1/5$、$2/5$、$3/5$、$7/5$、$8/5$ 等数值处被观察到。这些序列的共同要素是分母为奇数。然而，解释这些实验的原始困难并不在于说明霍尔电导为何取分数值，而在于阐明展现这种分数量化的电子态的本质。例如，一个最低朗道能级在填充 $\nu$ 下被分数占据的无相互作用模型将表现出电导 $\nu e^{2}/h$。然而，当一个朗道能级被分数占据时，电子可以散射进空态，因而纵向输运将不是无耗散的；结果 $\sigma_{xx}\neq 0$。此外，平台在分数填充处的持续存在表明，一个适当部分占据的朗道能级不知怎么对哪怕最低的激发也是稳定的。也就是说，有一个能隙把基态与所有激发态隔开。在不存在电子相互作用时，往部分填充的朗道能级中添加一个额外电子所需能量几乎为零。这意味着，在无相互作用模型中，随着磁场改变，部分填充朗道能级的电导陡峭平台应当消失。

因此，不把电子相互作用的作用包括进来，分数量子霍尔态便无法理解。然而，一开始我们并不清楚这个态应当是什么样子。至此，我们已经引入了两种本质上源于电子–电子相互作用的电子态：维格纳晶体与超导态。纵向电阻的消失暗示分数量子霍尔态与超导态有几分相似，但这样的态如何能在大的垂直磁场中存留下来，并不清楚。那么维格纳晶体呢？如第 5 章所述，由于磁场冻结了电子的零点运动，维格纳结晶得以稳定。事实上，在粒子间距相对于回旋半径（磁性长度）较大的极限下，电子关联占主导地位，维格纳结晶的条件变得有利。然而，维格纳晶体并不表现出无耗散输运，因为必须先施加一个阈值电压输运才会出现。于是我们猜想，分数量子霍尔态的答案另有所在。

确实另有所在。沿劳夫林（Laughlin）（文献 [L1983; L1987]）的思路，我们考虑处于垂直磁场中的相互作用电子气
\begin{equation*}
H=\sum_{i}\left[\frac{1}{2m}\left(\mathbf{p}_{i}-\frac{e}{c}\mathbf{A}_{i}\right)^{2}+V(\mathbf{r}_{i})\right]+\frac{1}{2}\sum_{\stackrel{i,j}{i\neq j}}\frac{e^{2}}{|\mathbf{r}_{i}-\mathbf{r}_{j}|},\tag{15.83}
\end{equation*}
其中 $V(\mathbf{r})$ 是离子提供的补偿性中和势，对 $i$ 与 $j$ 的求和遍及电子。霍尔电导的分数量化在 $15\,\mathrm{T}$ 的磁场下被观察到。在此场强下，回旋能量大约是库仑能量 $e^{2}/\ell$ 的三倍。因此，用最低朗道能级的无相互作用本征态作为构造真实多体波函数的出发点，应当是相当好的近似。一般地，无相互作用本征态是电子坐标的多项式与高斯型因子的乘积。对于相互作用系统，我们预期真实的多体态将涉及电子坐标之差。让我们定义复电子坐标 $z=x+\mathrm{i}y$。我们考虑如下拟设（ansatz）：
\begin{equation*}
\Psi_{N}(\mathbf{r}_{1},\ldots,\mathbf{r}_{N})=\prod_{1\le j<k\le N}f(z_{j}-z_{k})\exp\left(-\sum_{j=1}^{N}\frac{|z_{j}|^{2}}{4\ell^{2}}\right),\tag{15.84}
\end{equation*}
其中 $f(z_{j}-z_{k})$ 是电子坐标的多项式。在研究了由式 (15.83) 所描述的双电子问题之后，劳夫林发现 $f(z_{1}-z_{2})=(z_{1}-z_{2})^{2p+1}$，其中 $p$ 为整数。由于指数 $2p+1$ 为奇数，波函数在交换两个电子时是反对称的。对等价的三电子问题也得到了类似的结果（文献 [L1983]）。于是，劳夫林提出
\begin{equation*}
\Psi_{m}(\mathbf{r}_{1},\ldots,\mathbf{r}_{N})=\prod_{1\le j<k\le N}(z_{j}-z_{k})^{m}\exp\left(-\sum_{j=1}^{N}\frac{|z_{j}|^{2}}{4\ell^{2}}\right)\tag{15.85}
\end{equation*}
必然准确地描述分数量子霍尔系统的基态。虽然这个波函数最初被论证为一个变分态（其中 $m$ 待定），但对形如 $u(r)=1/r$、$-\ln r$ 与 $\exp(-r^{2}/2)$ 的相互作用势，该波函数与小型电子团簇的精确本征态之间的重叠通常大于 $99\%$。这表明 $\Psi_{m}$ 的变分成分极小，$m$ 必须由某条基本原理来确定。结果是，$\Psi_{m}$ 是总角动量算符的本征态，本征值为
\begin{equation*}
M=\frac{N(N-1)m}{2}.\tag{15.86}
\end{equation*}
$M$ 之所以是算符 $L_{z}$ 的总角动量，是因为把式 (15.84) 中 $z_{j}-z_{k}$ 的乘积展开后可以看到，每一项中 $z_{i}$ 的因子至多有 $N(N-1)m/2$ 个。若把 $L_{z}$ 算符的微分形式作用在这样的乘积上，便得到本征值 $M$。因此，我们可以把劳夫林态看作最低朗道能级内具有相同角动量的诸态的叠加。这完全消除了劳夫林态的变分成分。

为了准确理解劳夫林态究竟描述什么样的物理，我们把这个态的平方
\begin{equation*}
|\Psi_{m}|^{2}=\exp[-\beta\Phi]\tag{15.87}
\end{equation*}
写成玻尔兹曼权重，其中 $\Phi$ 是相互作用能，
\begin{equation*}
\Phi(z_{1},\ldots,z_{N})=-2\sum_{1\le j<k\le N}\ln|z_{j}-z_{k}|+\frac{1}{2m\ell^{2}}\sum_{j=1}^{N}|z_{j}|^{2}\tag{15.88}
\end{equation*}
而 $\beta=m$ 起逆温度的作用。式 (15.88) 恰好是由 $N$ 个电荷为 $\sqrt{2}$ 的相同电荷组成的单成分等离子体（one-component plasma）（例如电子气）的相互作用能。式 (15.88) 的第二项表示与中和背景的相互作用能 $U_{\mathrm{b}}(z_{i})=|z_{i}|^{2}/(2m\ell^{2})$。为清楚地看到这一点，我们注意到补偿背景的势应当满足泊松方程 $\nabla^{2}U_{\mathrm{b}}(z_{i})=4\pi\rho$，其中 $\rho$ 由 $\nu n_{\mathrm{B}}=\nu/(2\pi\ell^{2})$ 给出（见式 (15.21)）。在泊松方程中完成微商便知，单位面积的补偿电荷密度
\begin{equation*}
\rho=\frac{1}{2\pi\ell^{2}m}\equiv\frac{\nu}{2\pi\ell^{2}},\tag{15.89}
\end{equation*}
恰恰就是最低朗道能级中均匀的电子电荷密度，只要我们把朗道填充因子 $\nu$ 认定为 $1/m$。对于使最低朗道能级的填充为 $\nu=1/m$ 的密度与磁场强度，劳夫林态准确地描述 $N$ 电子系统的基态。作为具有均匀电子密度的系统，劳夫林态不同于电子晶体态，例如维格纳晶体。事实上，凝聚成劳夫林态的电子并不破坏任何对称性。因此，上一章给出的朗道型描述——它必然涉及某个序参量的对称性破缺——不可能用于分数量子霍尔态的形成。劳夫林态对分数量子霍尔效应发展的重要性，无论怎样评价都不为过。

我们现在蓄势待发，可以探究劳夫林液体的激发了。当填充分数 $\nu$ 不同于 $1/m$ 时，激发便会出现。考虑用一个无限细的磁性螺线管在 $z_{0}$ 处刺穿样品来改变填充。虽然此时 $\nu$ 略小于 $1/m$，电子仍会试图停留在态 $\Psi_{m}$ 中。然而，若不降低磁性螺线管插入点处的电荷密度，它们便无法做到这一点。我们可以通过把电子从 $z_{0}$ 处排除来模拟这种耗减。于是我们预期
\begin{align*}
\Psi_{m}^{+}(z_{0};z_{1},\ldots,z_{N})&=\prod_{j=1}^{N}(z_{j}-z_{0})\Psi_{m}(z_{1},\ldots,z_{N})\\
&=A_{z_{0}}^{+}\Psi_{m}(z_{1},\ldots,z_{N})\tag{15.90}
\end{align*}
可能描述在 $z_{0}$ 处带有“准空穴”（quasi-hole）的新多体态的波函数。正如劳夫林所证明的（文献 [L1983]），事实确实如此。对式 (15.90) 中准粒子波函数的一个更定量的论证来自如下观察：若螺线管携带磁通 $\phi$，则每个单粒子态相应地变为
\begin{equation*}
z^{n}\mathrm{e}^{-\frac{|z|^{2}}{4\ell^{2}}}\to z^{n+\alpha}\mathrm{e}^{-\frac{|z|^{2}}{4\ell^{2}}},\tag{15.91}
\end{equation*}
以容纳附加的磁通。这里 $\alpha=\phi/\phi_{0}$，$\phi_{0}=hc/e$ 为磁通量子。若螺线管携带一个磁通量子，则高斯型因子的前置系数为 $z^{n+1}$。每个单粒子态现在都额外多乘了一个 $z$ 因子，这为式 (15.90) 的准粒子波函数拟设提供了支持。为了利用等离子体类比，我们把准粒子波函数平方，发现除一个背景归一化因子 $|z_{0}|^{2}$ 之外，带有“准空穴”的多体态的能量
\begin{equation*}
\Phi_{qp}(z_{0};z_{1},\ldots,z_{N})=\Phi(z_{1},\ldots,z_{N})-\frac{2}{m}\sum_{j=1}^{N}\ln|z_{j}-z_{0}|,\tag{15.92}
\end{equation*}
正是单成分等离子体与一个固定在 $z_{0}$ 处的电荷相互作用的能量。该电荷的大小为 $1/m$，或者以电子电荷为单位，$e/m$。同样，若螺线管抽走一个磁通量子，则会产生一个电荷为 $-e/m$ 的“准电子”。数值上，在三电子分数量子霍尔态中产生或湮灭准粒子的能量，在 $15\,\mathrm{T}$ 磁场下约为 $4\,\mathrm{K}$（文献 [L1983]）。分数量子霍尔态能隙的改进估算由 Girvin、MacDonald 与 Platzman 给出（文献 [GMP1985]）。事实上，他们的工作对确立“劳夫林态的所有集体激发都具有有限能隙”起了关键作用。因此，劳夫林态确实满足对所有激发都有能隙这一判据。激发能可以看作把一个电荷为 $\pm e/m$ 的粒子放进量子液体所需的库仑能。电荷作用的距离正比于磁性长度，即 $\propto\ell$。因此，激发能应当按 $e^{2}/\ell\propto\sqrt{B}$ 标度，从而在无外加磁场时消失。正是激发谱中能隙的存在使劳夫林态不可压缩，并最终导致无耗散输运。

尽管如此，一些实验在二维量子霍尔系统中揭示了无能隙激发的存在（文献 [ASH1983; MDF1985]）。这些激发被认为居于量子霍尔系统的边缘，因为劳夫林的规范论证明确地把它们从体内排除。对于整数量子霍尔态，边缘激发可以很好地用费米液体理论描述，因为在整数效应中电子相互作用相对不重要。然而，正如我们所看到的，在分数量子霍尔效应中情况完全不同。局限于分数量子霍尔系统边缘运动的量子力学态，本质上是一维强关联电子系统。我们在第 10 章中已经证明，一维系统中的电子相互作用导致卢廷格液体而非费米液体行为。由于边缘电流的手性，分数量子霍尔效应中的边缘态是手性卢廷格液体（文献 [W1990]）。它们表现出第 9 章所讨论的卢廷格液体的全部标志性性质，其中包括在费米能附近按代数方式趋于零的激发谱。图 10.6 证实了 $\nu=1/3$ 量子霍尔态边缘激发谱的代数依赖关系，从而使分数量子霍尔效应边缘态的手性卢廷格液体模型立于坚实的实验基础之上。

现在暂且考虑与交换分数量子霍尔态中的两个准空穴或两个准电子相联系的统计问题（文献 [A1985; F1992]）。描述这种成对激发的波函数与式 (15.90) 类似，只是现在乘在劳夫林态 $\Psi_{m}$ 上的是乘积 $A_{z'}^{\pm}A_{z''}^{\pm}$，我们把两个准粒子定位于 $z=z'$ 与 $z=z''$。我们要讨论的问题是：在两个准粒子的互换使得
\begin{equation*}
A_{z'}^{\pm}A_{z''}^{\pm}\Psi_{m}(z_{1},\ldots,z_{N})=\mathrm{e}^{i\phi}A_{z''}^{\pm}A_{z'}^{\pm}\Psi_{m}(z_{1},\ldots,z_{N}),\tag{15.93}
\end{equation*}
时，波函数获得的相位 $\phi$ 是多少？对于电子的交换，$\phi=\pi$；对于玻色子，$\phi=2\pi$，或等价地 $0$。我们现在就将证明，劳夫林态中两个准粒子互换的相位变化是 $1/m$，因而是分数的。

相当一般地，在矢量势存在时沿闭合回路运动的粒子，其波函数会获得相位
\begin{equation*}
i\gamma=i\frac{q}{\hbar c}\oint\mathrm{d}\ell\cdot\mathbf{A}.\tag{15.94}
\end{equation*}
我们将把矢量势取为粒子所感受的磁场 $\mathbf{B}$ 的生成者。因此，决定相位的线积分就简单地等于场乘以路径所围的面积。设回路的半径为 $R$，则相位变化为 $\gamma=q\pi R^{2}B/\hbar c$。回忆磁场通过 $\rho=\nu eB/\hbar c$ 与电子密度相联系，而准粒子的电荷为 $e\nu$，我们发现总共经历的相位为 $\gamma=2\pi N$，其中 $N$ 是系统中的电子数。这是一个关键结果，因为相位正比于回路所围电子的数目。现在让我们重做上述论证，假定电子海之中有一个电荷为 $e/m$ 的准空穴。此时我们必须考虑准空穴导致的密度耗减。对一个准空穴，耗减由填充 $\nu$ 或 $1/m$ 给出。于是相位现在由 $2\pi(N-1/m)$ 给出。设 $|z'|<|z''|$。这样，当 $z''$ 处的准空穴沿一个环绕 $z'$ 处准空穴的回路运动时，相位变化为 $2\pi(N-1/m)$。$2\pi N$ 因子不改变波函数的符号。因此，额外获得的相位是 $2\pi/m$。然而，要互换两个准粒子，只需走过半圈，再把两个粒子平移。由于平移不产生任何相位变化，两个准粒子互换的净相位变化是 $\pi/m$。对 $m=1$，我们回到电子的结果。然而，对任何分数填充，相位变化都是分数。因此，劳夫林液体中的准粒子既表现出分数电荷，又表现出分数统计。遵守分数统计的粒子（文献 [W1982]）被称为任意子（anyon），因为它们在互换时可以获得任意的相位变化。

在本节收尾之际，我们引入对劳夫林态的一个看似无害的改写（文献 [J1989]）
\begin{equation*}
\Psi_{m}(\mathbf{r}_{1},\ldots,\mathbf{r}_{N})=\prod_{1\le j<k\le N}(z_{j}-z_{k})^{m-1}\Psi_{1}(\mathbf{r}_{1},\ldots,\mathbf{r}_{N})\tag{15.95}
\end{equation*}
这里 $\Psi_{1}(\mathbf{r}_{1},\ldots,\mathbf{r}_{N})$ 就是 $m=1$ 的劳夫林态，因而对应于完全填满的朗道能级的波函数。此式中的前置系数类似于“准空穴”波函数中附着通量的前置系数。劳夫林态可以看作一个完全填满的朗道能级，其中每个电子都实实在在地附着了 $m-1$ 个磁通量子。虽然电子绝不可能像书页装订成书那样真正束缚磁通，但通量附着方案为思考分数量子霍尔效应提供了一种图像化的方式。此外，它还有这样一个优点：分数量子霍尔效应可以被看成整数量子霍尔效应的一个变种，只是其中的虚拟粒子携带着电荷 $e$ 与 $m-1$ 个单位的磁通。作为一种计算工具，通量附着方案或复合粒子图像在生成最低朗道能级中分数态的层级结构方面极为有用，因而构成了当前量子霍尔问题各种场论方法的基础（文献 [F1991; LF1991; HLR1993; KLZ1992]）。事实上，在偶分母态的语境下，复合费米子方法已经扮演了中心角色。$\nu=1/2$ 的偶分母态特别奇特（文献 [W1993]）：它看起来并不是一个具有消失的纵向电阻的真分数态。$\nu=1/2$ 处的通量附着方案导致一种复合费米子液体，其中复合费米子感受到的平均磁场为零。在这样的系统中，既然磁场表观上消失，就应当涌现出一个费米面。关于这一问题的近期工作（文献 [R1998; SM1997]）表明，$\nu=1/2$ 态是一个可压缩态，由具有偶极矩的费米型激发组成。

当前感兴趣的另一个问题是量子霍尔效应的低磁场极限。当垂直磁场足够低、使填充分数超过 $\nu>7/2$ 时，实验观察到（文献 [D1999; L1999]）半整数态附近的纵向电阻率依赖于电流的方向，即 $\rho_{xx}\neq\rho_{yy}$。各向异性的输运与均匀电子态不相容。原则上，当磁场足够低以致任何分数量子霍尔态都不可及、但又不低到使无序破坏由回旋频率 $\omega_{\mathrm{c}}$ 所建立的能隙时，二维电子气的基态预期是非均匀的。在这个中间磁场范围内，基态由短程交换相互作用与长程直接库仑相互作用的竞争所支配。由于这两种相互作用符号相反，电子气破裂成一个个畴（domain）（文献 [FPA1979; FKS1996; MC1996; RHY1999; J1999; SMP2000]）。当前的理论工作致力于表述有效低能理论（文献 [FK1999; MF2000; BF2001; B2001; L2001]），以及低场极限下存在电荷调制时的输运性质。

\section*{小结}
\addcontentsline{toc}{section}{小结}

本章我们给出了量子霍尔效应背后的基本物理。整数效应源于如下边界条件：当整数个朗道能级被填满时，载流态或边缘态必须单值。在分数量子霍尔效应中，部分占据最低朗道能级的电子凝聚成一个强关联态（劳夫林态），其能量低于相应的维格纳晶体。这个态是不可压缩的，即它对所有激发都有能隙，并表现出电荷为 $\pm e/m$ 的分数电荷激发。这类带电激发的激发能按 $1/\ell$ 标度，因而在零场下消失。正是由于这个原因，虽然劳夫林态由强电子关联所中介，但其不可压缩性来源于大的垂直磁场。当前感兴趣的量子霍尔问题包括 $\nu=1/2$ 处表观上类费米液体的态，以及磁场足够低、填充超过 $\nu=7/2$ 时各向异性态的出现。

\section*{习题}
\addcontentsline{toc}{section}{习题}

\textbf{习题 15.1}\quad 考虑二维哈密顿量
\begin{equation*}
p_{x}\sigma^{x}+p_{y}\sigma^{y}+m(x)\sigma^{z},\tag{15.96}
\end{equation*}
其中
\begin{equation*}
m(x)=\begin{cases}m, & x>0,\\ -m, & x<0,\end{cases}
\end{equation*}
（a）证明形如
\begin{equation*}
\psi(x,y)=\mathrm{e}^{ip_{y}y}\phi(x)\tag{15.97}
\end{equation*}
而
\begin{equation*}
\phi(x)=\exp\left[-\int_{0}^{x}m(x')\,\mathrm{d}x'\right]\phi_{0},\tag{15.98}
\end{equation*}
其中 $\phi_{0}$ 为二分量旋量的波函数，描述局域在边界 $x=0$ 处的零能态。也就是说，请显式地求出 $\phi_{0}$。（b）再证明同一个态也描述一个束缚在 $x=0$ 处但有色散的态。这个态的能量是多少？为什么这样的态只有一个？

\textbf{习题 15.2}\quad （a）一个完全局限于 $x$–$y$ 平面内运动的单粒子由单粒子波函数
\begin{equation*}
\Psi_{m}(z)=z^{m}\exp\left(\frac{-|z|^{2}}{4\ell^{2}}\right),\tag{15.99}
\end{equation*}
描述，其中 $z=x+\mathrm{i}y$。证明该波函数是 $L_{z}$ 的本征态。（b）现在考虑推广到局限于 $x$–$y$ 平面内运动的 $N$ 个粒子：
\begin{equation*}
\Psi_{m}(r_{1},\ldots,r_{N})=\prod_{j<k}(z_{j}-z_{k})^{m}\exp\left(-\sum_{j=1}^{N}\frac{|z_{j}|^{2}}{4\ell^{2}}\right),\tag{15.100}
\end{equation*}
其中 $z_{j}=x_{j}+\mathrm{i}y_{j}$。对哪些 $m$ 值，此波函数在粒子互换下是反对称的？（c）证明该波函数是总 $L_{z}$ 算符的本征态，本征值为
\begin{equation*}
M=\frac{N(N-1)m}{2}.\tag{15.101}
\end{equation*}
对 $N$ 个粒子，总 $L_{z}$ 算符就是每个粒子的 $L_{z}$ 之和。

\textbf{习题 15.3}\quad 证明式 (15.29) 是式 (15.48) 的时间反演算符。

\textbf{习题 15.4}\quad 由 $T^{2}=-1$（$T$ 为时间反演算符）这一事实，证明克拉默斯倍制定理。

\section*{参考文献}
\addcontentsline{toc}{section}{参考文献}

\begin{itemize}[leftmargin=2em,itemsep=0pt,parsep=0pt,label={}]
\item [ASH1983] S. J. Allen, Jr., H. L. Störmer, and J. C. M. Hwang, \emph{Phys. Rev. B} \textbf{28}, 4875 (1983).
\item [AMY1975] T. Ando, Y. Matsumoto, and Y. Uemura, \emph{J. Phys. Soc. Jpn.} \textbf{39}, 279 (1975).
\item [AA1981] H. Aoki and T. Ando, \emph{Solid St. Comm.} \textbf{38}, 1079 (1981).
\item [A1985] D. P. Arovas, in \emph{Geometric Phases in Physics}, eds. A. Shapere and F. Wilczek (World Scientific, Singapore, 1985), p. 284.
\item [BF2001] D. G. Barci and E. Fradkin, cond-mat/0106171.
\item [B2001] D. G. Barci, E. Fradkin, S. A. Kivelson, and V. Oganeysan, cond-mat/0105448.
\item [BHZ2006] B. A. Bernevig, T. L. Hughes, and S.-C. Zhang, \emph{Science} \textbf{314}, 1757 (2006).
\item [BZ2006] B. A. Bernevig and S.-C. Zhang, \emph{Phys. Rev. Lett.} \textbf{96}, 106802 (2006).
\item [CC1988] J. T. Chalker and P. D. Coddington, \emph{J. Phys. C} \textbf{21}, 2665 (1988).
\item [D1999] R. R. Du, D. C. Tsui, H. L. Störmer, L. N. Pfeiffer, K. W. Baldwin, and K. W. West, \emph{Solid St. Commun.} \textbf{109}, 389 (1999).
\item [FKS1996] M. M. Fogler, A. A. Koulakov, and B. I. Shklovskii, \emph{Phys. Rev. B} \textbf{54}, 1853 (1996).
\item [F1992] S. Forte, \emph{Rev. Mod. Phys.} \textbf{64}, 193 (1992).
\item [F1991] E. Fradkin, \emph{Field Theories of Condensed Matter Systems}, (Addison-Wesley, Redwood City, USA, 1991).
\item [FK1999] E. Fradkin and S. Kivelson, \emph{Phys. Rev. B} \textbf{59}, 8065 (1999).
\item [FPA1979] H. Fukuyama, P. Platzman, and P. W. Anderson, \emph{Phys. Rev. B} \textbf{19}, 5211 (1979).
\item [GMP1985] S. M. Girvin, A. H. MacDonald, and P. M. Platzman, \emph{Phys. Rev. Lett.} \textbf{54}, 581 (1985).
\item [H1879] E. H. Hall, \emph{Am. J. Math.} \textbf{2}, 287 (1879).
\item [H1982] B. I. Halperin, \emph{Phys. Rev. B} \textbf{25}, 2185 (1982).
\item [HLR1993] B. I. Halperin, P. A. Lee, and N. Read, \emph{Phys. Rev. B} \textbf{47}, 7312 (1993).
\item [HK2010] M. Z. Hasan and C. L. Kane, \emph{Rev. Mod. Phys.} \textbf{82}, 3045 (2010).
\item [H1937] H. Hellmann, \emph{Einführung in die Quantenchemie} (Franz Denticke, Leipzig, Germany, 1937) p. 285; see also R. P. Feynman, \emph{Phys. Rev.} \textbf{56}, 340 (1939).
\item [H2008] D. Hsieh, D. Qian, L. Wray, Y. Xia, Y. S. Hor, R. J. Cava, and M. Z. Hasan, \emph{Nature} \textbf{452}, 970 (2008).
\item [H2009] D. Hsieh, Y. Xia, D. Qian, \emph{et al.}, \emph{Nature} \textbf{460}, 1101 (2009).
\item [J1989] J. Jain, \emph{Phys. Rev. B} \textbf{40}, 8079 (1989).
\item [J1999] T. Jungwirth, A. H. MacDonald, L. Smrcka, amd S. M. Girvin, \emph{Phys. Rev.} \textbf{60}, 15574 (1999).
\item [KM2005a] C. L. Kane and E. J. Mele, \emph{Phys. Rev. Lett.} \textbf{95}, 226801 (2005).
\item [KM2005b] C. L. Kane and E. J. Mele, \emph{Phys. Rev. Lett.} \textbf{95}, 146802 (2005).
\item [KLZ1992] S. Kivelson, D.-H. Lee, and S.-C. Zhang, \emph{Phys. Rev. B} \textbf{46}, 2223 (1992).
\item [KDP1980] K. von Klitzing, G. Dorda, and M. Pepper, \emph{Phys. Rev. Lett.} \textbf{45}, 494 (1980).
\item [K2007] M. König, S. Wiedmann, C. Brüne, \emph{et al.}, \emph{Science} \textbf{318}, 766 (2007).
\item [L1980] R. B. Laughlin, \emph{Phys. Rev. B} \textbf{23}, 5632 (1980).
\item [L1983] R. B. Laughlin, \emph{Phys. Rev. B} \textbf{27}, 3383 (1983); \emph{Phys. Rev. Lett.} \textbf{50}, 873 (1983).
\item [L1987] R. B. Laughlin, in \emph{The Quantum Hall Effect}, eds. E. Prange and S. M. Girvin (Springer, Berlin, 1987) p. 233.
\item [LWK1993] D.-H. Lee, Z. Wang, and S. A. Kivelson, \emph{Phys. Rev. Lett.} \textbf{70}, 4130 (1993).
\item [LC1994] D. K. K. Lee and J. T. Chalker, \emph{Phys. Rev. Lett.} \textbf{72}, 1510 (1994).
\item [L1999] M. P. Lilly, K. B. Cooper, J. P. Eisenstein, L. N. Pfeiffer, and K. W. West, \emph{Phys. Rev. Lett.} \textbf{82}, 394 (1999).
\item [L2001] A. Lopatnikova, S. H. Simon, B. I. Halperin, and X.-G. Wen, cond-mat/0105079.
\item [LF1991] A. Lopez and E. Fradkin, \emph{Phys. Rev. B} \textbf{44}, 5246 (1991).
\item [MF2000] A. H. MacDonald and M. P. A. Fisher, \emph{Phys. Rev. B} \textbf{61}, 5724 (2000).
\item [MT1999] J. B. Marston and S.-W. Tsai, \emph{Phys. Rev. Lett.} \textbf{82}, 4906 (1999).
\item [MDF1985] D. B. Mast, A. J. Dahm, and A. L. Fetter, \emph{Phys. Rev. Lett.} \textbf{54}, 1706 (1985).
\item [MC1996] R. Moessner and J. T. Chalker, \emph{Phys. Rev. B} \textbf{54}, 5006 (1996).
\item [MB2007] J. E. Moore and L. Balents, \emph{Phys. Rev. B} \textbf{75}, 121306 (R) (2007).
\item [P1981] R. E. Prange, \emph{Phys. B} \textbf{23}, 4802 (1981).
\item [P1987] R. E. Prange, in \emph{The Quantum Hall Effect}, eds. R. E. Prange and S. M. Girvin (Springer, Berlin, 1987) p. 1.
\item [P1984] A. Pruisken, \emph{Nucl. Phys. B} \textbf{235} (FS11), 277 (1984).
\item [QHZ2008] X.-L. Qi, T. L. Hughes, and S.-C. Zhang, \emph{Phys. Rev. B} \textbf{78}, 195424 (2008).
\item [QZ2010] X.-L. Qi and S.-C. Zhang, \emph{Phys. Today} \textbf{63}, 33 (2010).
\item [R1998] N. Read, \emph{Phys. Rev. B} \textbf{58}, 16262 (1998).
\item [RHY1999] E. H. Rezayi, F. D. M. Haldane, and K. Yang, \emph{Phys. Rev. Lett.} \textbf{83}, 1219 (1999).
\item [R2009] R. Roy, \emph{Phys. Rev. B} \textbf{79}, 195321 (2009).
\item [SM1997] R. Shankar and G. Murthy, \emph{Phys. Rev. Lett.} \textbf{79}, 4437 (1997).
\item [SMP2000] T. D. Stanescu, I. Martin, and P. Phillips, \emph{Phys. Rev. Lett.} \textbf{84}, 1288 (2000).
\item [TFK2008] J. C. Y. Teo, L. Fu, and C. L. Kane, \emph{Phys. Rev. B} \textbf{78}, 045426 (2008).
\item [T1982] D. Thouless, M. Kohmoto, M. denNijs, and M. Nightingale, \emph{Phys. Rev. Lett.} \textbf{49}, 405 (1982).
\item [T1983] S. A. Trugman, \emph{Phys. Rev. B} \textbf{27}, 7539 (1983).
\item [TSG1982] D. C. Tsui, H. L. Störmer, and A. C. Gossard, \emph{Phys. Rev. Lett.} \textbf{48}, 1599 (1982).
\item [W1990] X. G. Wen, \emph{Phys. Rev. B} \textbf{41}, 12838 (1990).
\item [W1982] F. Wilczek, \emph{Phys. Rev. Lett.} \textbf{49}, 957 (1982).
\item [W1993] R. L. Willett, R. R. Ruel, M. A. Paalanen, K. W. West, and L. N. Pfeiffer, \emph{Phys. Rev. B} \textbf{47}, 7344 (1993).
\item [X2009] Y. Xia, D. Qian, D. Hsieh, \emph{et al.}, \emph{Nature Physics} \textbf{5}, 398 (2009).
\item [Z2009] H. Zhang, C.-X. Liu, X.-L. Qi, X. Dai, Z. Fang, and S.-C. Zhang, \emph{Nature Physics} \textbf{5}, 438 (2009).
\item [Z1999] M. R. Zirnbauer, hep-th/9905054.
\end{itemize}
